\section{IAM - What should you know by now}
\begin{itemize}
    \item Users: Long term credentials
    \item Groups
    \item Roles: short-term credentials, uses STS
    \item EC2 Instance Roles: Uses the EC2 metadata service. One role of a time per instance
    \item Service Roles: API Gateways, CodeDeploy etc
    \item Cross Account roles
    \item Policies
    \item AWS Managed
    \item Customer Managed
    \item Inline Policies
    \item Resource Based Policies (S3 Bucket, SQS queues, etc......)
\end{itemize}

\subsection{IAM Policies Deep Dive}
\begin{itemize}
    \item Anatomy of a policy: JSON doc with Effect, Action, Resource, Conditions, Policy Variables
    \item Explicit DENY has precedence over ALLOW
    \item Best practice: use least privilege for maximum security
    \item Access Advisor: See permissions granted and when last accessed
    \item Access Analyser: Analyse resources that are shared with external entity
    \item Navigate Examples at: %put link to AWS docs here
\end{itemize}

\subsection{IAM AWS Managed Policies - Administrator Access example}

\begin{lstlisting}[language=JSON, caption=Sample JSON Data, label=lst:json-example]
{
"Version": "2012-10-17",
"Statement": [
{
"Effect": "Allow",
"Action": "*",
"Resource": "*"
}
]
}
\end{lstlisting}
% Add other listings as required

\subsection{IAM Policies Conditions}
% Insert IAM policies conditions

\subsection{IAM Policies Variables and Tags}
% Insert listings here

\subsection{IAM Roles vs Resource Based Policies}
\begin{itemize}
    \item Attach a policy to a resource (example: S3 bucket policy) versus attaching of a using a role as a proxy
%insert a diagram
    \item When you assume a role (user, application or service), you give up your original permissions and take the permissions assigned to the role.
    \item When using a resource-based policy, the principal doesn't have to give up any permissions
    \item Example: User in account A needs to scan a DynamoDB table in Account A and dump it in an S3 bucket in Account B
\end{itemize}

\subsection{IAM Permissions Boundaries}
\begin{itemize}
    \item IAM permission boundaries are supported for users and roles (not groups)
    \item Advanced feature to use a managed policy to set the maximum permissions an IAM entity can get
    \item Can be used in combinations of AWS organisations SCP
\end{itemize}

\subsection{Use cases}

\begin{itemize}
    \item Delegate responsibilities to non administrators within their permission boundaries, for example create create new IAM users
    \item Allow developers to self-assign policies and managed their own permissions while making sure they can't 'escalate' their privileges (i.e make themselves admin)
    \item Useful to restrict one specific user (instead of a whole account using Organisations and SCP)
\end{itemize}

\subsection{IAM Access Analyser}
\begin{itemize}
    \item Find out which resources are shared externally
    \item S3 Buckets
    \item IAM Roles
    \item KMS Keys
    \item Lambda Functions and Layers
    \item SQS queues
    \item Secrets Manager Secrets
    \item Define Zone ofTrust = AWS Account or AWS Organisation
    \item Access outside of zone trusts => findings
\end{itemize}

\subsection{IAM Access Analyser}
\begin{itemize}
    \item IAM Access Analyser Policy Validation
    \item Validates your policy against IAM policy grammar and best practices
    \item General warnings, security warnings errors suggestions
    \item Provides actionable recommendations
    \item IAM Access Analyser Policy Generation
    \item Generates IAM policy based on access activity
    \item CloudTrail logs is reviewed to generate the policy with the fine-grained permissions and the appropriate Actions and Services
    \item Reviews CloudTrail logs for up to 90 days
\end{itemize}


\section{STS}

\subsection{Using STS to Assume a Role}
\begin{itemize}
    \item Define an IAM Role within your account or cross-account
    \item Define which principals can access this IAM role
    \item Using AWS STS (Security Token Service) to retrieve credentials and impersonate the IAM role you have access to (AssumeRole API)
    \item Temporary credentials can be valid between 15 to 12 hours.
\end{itemize}

\subsection{Assuming a Role with STS}
\begin{itemize}
    \item Provide access for an IAM user in one AWS account that you own to access resources in another account that you own.
    \item Provide access to IAM users in AWS accounts owned by third parties
    \item Provide access for services offered by AWS to AWS resources
    \item Provide access for externally authenticated users (identity federation)
    \item Ability to revoke active sessions and credentials for a role (by adding a policy using a time statement - AWSRevokeOlderSessions)
\end{itemize}

When you assume a role (user, application or service), you give up your original permissions and take the permissions assigned to the role.

\subsection{Providing Access to an IAM User in Yours or Another AWS Account That You Own}
\begin{itemize}
    \item You can grant your IAM users permissions to switch to roles within your AWS account or to roles defined in other AWS accounts that you own.
%Diagram
    \item Benefits:
    \item You must explicitly grant your users permission to assume the role
    \item Your users must actively switch to the role using the AWS management console or assume the role using the AWS CLI or AWS API
    \item You can add multi-factor authentication (MFA) protection to the role so that only users who sign in with an MFA device can assume the role.
    \item Least Privilege + auditing using CloudTrail
%Diagram
\end{itemize}

\subsection{Providing Access to AWS Accounts Owned by Third Parties}
\begin{itemize}
    \item Zone of trust = accounts, organisations that you own
    \item Outside Zone ofTrust = 3rd parties
    \item Use IAM Access Analyser to find out which resources are exposed
    \item For granting access to a thrid party:
    \item The third part AWS account ID
    \item An External ID (secret between you and the third party)
    \item To uniquely associate with the role between you and 3rd party
    \item Must be provided when defining the trust and when assuming the role
    \item Must be chosen by the third party
    \item Define permissions in the IAM policy
\end{itemize}

%The confused deputy, diagram

\subsubsection{Session Tags in STS}
\begin{itemize}
    \item Tags that you pass when you assume an IAM Role or federate user in STS
    \item aws:PrincipalTag Condition
    \item Compares the tags attached to the principal making the request with the tag you specified in the policy.
    \item Example: Allow a principal to pass session tags only if the principal making the reqeust has the specified tags.
% diagram
\end{itemize}

\subsubsection{STS Important APIs}
\begin{itemize}
    \item AssumeRole: access a role within your account or cross-account
    \item AssumeRoleWithSAML: return credentials for users logged with SAML
    \item AssumeRoleWithWebIdentity: return creds for users logged with an IdP
    \item Example providers include Amazon Cognito, Login with Amazon, Facebook, Google or any OpenID Connect-compatible identity provider.
    \item AWS Recommends using Cognito
    \item GetSessionToken: for MFA from a user or AWS account root user
    \item GetFederationToken: obtain temporary creds for a federated user, usually a proxy app that will give the creds for a federated user, usually a proxy app that will give the creds to a distributed app inside a corporate network.
\end{itemize}

\subsection{Identity Federation in AWS}
\begin{itemize}
    \item Give users outside of AWS permissions to access AWS resources in your account
    \item You don't need to create IAM Users (user management is outside AWS)
    \item Use cases:
    \begin{itemize}
        \item A corporate has its own identity system (e.g Active Directory)
        \item Web / Mobile application that needs access to AWS resources
        \item Identity Federation can have many flavours:
        \item SAML 2.0
        \item Custom Identity Broker
        \item Web Identity Federation With(out) Amazon Cognito
        \item IAM Identity Centre
    \end{itemize}
\end{itemize}

\subsubsection{SAML 2.0 Federation}
\begin{itemize}
    \item Security Assertion Markup Language 2.0 (SAML 2.0)
    \item Open standard used by many identity providers (e.g ADFS)
    \item Supports integration with Microsoft Active Directory Federations Services (ADFS)
    \item Or any SAML 2.0 - compatible IdPs with AWS
    \item Access the AWS Console, AWS CLI or AWS API using temporary credentials
    \item No need to create IAM Users for each of your employees
    \item Need to setup a trust between AWS IAM and SAML 2.0 Identity Provider (both ways)
    \item Under-the-hood: Uses the STS API AssumeRoleWithSAML
    \item SAML 2.0 Federation is the "old way", IAM Identity Center Federation is the new managed and simpler way
%Insert a URL
\end{itemize}

\subsubsection{SAML 2.0 Federation - AWS API Access}
% Diagram here

\subsubsection{SAML 2.0 Federation - AWS Console Access}
% Diagram here

\subsubsection{SAML 2.0 Federation - Active Directory DS (ADFS)}
% Diagram here

\subsubsection{Custom Identity Broker Application}
\begin{itemize}
    \item Use only if identity provider is NOT compatible with SAML 2.0
    \item The identity broker is NOT compatible with SAML 2.0
    \item The identity broker must determine the appropriate IAM Role
    \item Uses the STS API AssumeRole or GetFederationToken
% Diagrams
\end{itemize}

\subsubsection{Web Identity Federation - Without Cognito}
- Not recommended by AWS - use cognito instead
% Insert a diagram

\subsubsection{Web Identity Federation - With Cognito}
\begin{itemize}
    \item Preferred over for Web Identity Federation
    \item Create IAM Roles using Cognito with the least privilege needed
    \item Build trust between the OIDC IdP and AWS
    \item Cognito benefits:
    \item Supports anonymous users
    \item Supports MFA
    \item Data Synchronisation
    \item Cognito replaces a Token Vending Machine (TVM)
%Insert a diagram
\end{itemize}

\subsubsection{Web Identity Federation - IAM Policy}
\begin{itemize}
    \item After being authenticated with Web Identity Federation, you can identify the user with an IAM policy variable
    \item Example:
%Insert examples
\end{itemize}

\subsection{AWS Directory Services (AD)?}
\begin{itemize}
    \item Found on any Windows Server with AD Domain Services
    \item Database of objects: User Accounts, Computers, Printers, File Shares, Security Groups
    \item Centralised security management, create account, assign permissions
    \item Objects are organised in trees
    \item A group of trees
%diagrams
\end{itemize}

\subsubsection{What is ADFS (AS Federation Services)?}
\begin{itemize}
    \item ADFS provides Single Sign-On across applications
    \item SAML across 3rd party: AWS Console, Dropbox, Officer365, etc.....
%diagrams
\end{itemize}

\subsubsection{AWS Directory Services}
\begin{itemize}
    \item AWS Managed Microsoft AD
    \item Create your own AD in AWS, manage users, locally supports MFA
    \item Establish "trust" connection with your own on premises
%Diagrams
    \item AD Connector
    \item Directory Gateway (proxy) to redirect to on-premises AD, supports MFA
    \item Users are managed on the on-premise AD
%Diagrams
    \item Simple AD
    \item AD-compatible managed directory on AWS
    \item Cannot be joined with on-premise AD
%Diagrams
\end{itemize}

\subsubsection{AWS Directory Services and AWS Managed Microsoft AD}
\begin{itemize}
    \item Managed Service: Microsoft AD in your AWS VPC
    \item EC2 Windows Instances:
    \item EC2 Windows instances can join the domain and run traditional AD applications (sharepoint, etc)
    \item Seamlessly Domain Join Amazon EC2 Instances from Multiple Accounts and VPCs
    \item Integrations
    \item RDS for SQL Server, AWS Workspaces, Quicksight.......
    \item AWS SSO to provide access to third party applications
    \item Standalone repository in AWS or jointed to on-premises AD
    \item Multi AZ deployment of AD in 2 AZ, # of DC (Domain Controllers) can be increased for scaling
    \item Automated backups
    \item Automated Multi-Region replication of your directory
% Insert Diagram
\end{itemize}

\subsubsection{AWS Microsoft Managed AD - Integrations}
%Add diagram

\subsubsection{Connect to on-premise AD}
\begin{itemize}
    \item Ability to connect your on-premise Active Directory to AWS Managed Microsoft AD
    \item Must establish a Direct Connection (DX) or VPN connection
    \item Can setup three kinds of forest trust
    \item One-way trust: AWS -> On-Premise
    \item One-way trust: On-Premise -> AWS
    \item Two-way forest: trust - AWS -> On-Premise
    \item Forest trust is different than synchronisation (replication is not supported)
\end{itemize}

\subsubsection{Solution Architecture: Active Directory Replication}
\begin{itemize}
    \item You may want to create a replica of your AD on EC2 in the cloud to minimise latency of in case DX or VPN goes down
    \item Establish trust between the AWS Manged Microsoft AD and EC2
% Diagram here
\end{itemize}

\subsubsection{AWS Directory Services AD Connector}
\begin{itemize}
    \item AD Connector is a directory gateway to redirect director requests to your on premises Microsoft Active Directory
    \item No caching capability
    \item Managed users solely on-premise, no possibility of setting up a trust
    \item VPN or Direct Connect
    \item Doesn't work with SQL Server, doesn't do seamless joining, can't share director.
%Insert diagram
\end{itemize}

\subsubsection{AWS Directory Services Simple AD}
\begin{itemize}
    \item Simple AD is an inexpensive Active Directory-compatible service with the common directory features
    \item Supports joining EC2 instances, manage users and groups
    \item Does not support MFA, RDS SQL server, AWS SSO
    \item Small: 500 users, large: 5000 users
    \item Powered by Samba 4, compatible with Microsoft AD
    \item lower cost, low scale, basic AD compatible or LDAP compatibility
    \item No trust relationship
\end{itemize}

\subsection{AWS Orgnisations}
%Diagram

\subsubsection{AWS Organisations - OrgnizationAccountAccessRole}
%Diagram
\begin{itemize}
    \item IAM role which grants full administrator permissions in the Member account to the Management account
    \item Used to perform admin tasks in the Member accounts (e.g - creating IAM users)
    \item Could be assumed by IAM users in the Management account
    \item Automatically added to all new Member account created with AWS organisations
    \item Must be created manually if you invite an existing Member account
\end{itemize}

\subsubsection{Multi Account Strategies}
\begin{itemize}
    \item Create accounts per department, per cost centre, per dev / test / prod, based on regulatory restrictions (using SCP), for better resource isolation (ex:VPC), to have separate per-account service limits isolated account for logging.
    \item Multi Account vs. On Account MultiVPC
    \item Use tagging standards for billing purposes
    \item Enable CloudTrail on all accounts, send logs to central S3 account
    \item Send CloudWatch logs to central logging account
    \item Strategy to create an account for security
\end{itemize}

\subsubsection{Orgnisational Units (OU) - Examples}
%Insert diagrams here

\subsubsection{AWS Organisation - Feature Modes}
\begin{itemize}
    \item Consolidated billing features:
    \item Consolidated Billing across all accounts - single payment method
    \item Pricing benefits from aggregated usage (volume discount for EC2, S3.....)
\end{itemize}

\subsubsection{All Features (Default)}
\begin{itemize}
    \item Includes consolidated billing features, SCP
    \item Invited accounts must approve enabling all features
    \item Ability to apply an SCP to prevent member accounts from leaving the org
    \item Can't switch back to Consolidated Billing Features only
\end{itemize}

\subsubsection{AWS Organisations - Reserved Instances}
\begin{itemize}
    \item For billing purposes, the consolidated billing features of AWS organisations treats all the accounts in the organisation as one account.
    \item The means that all accounts in the organisation can receive the hourly cost benefit of Reserved Instances that are purchased by any other account.
    \item The payer account (Management account) of an organisation can turn off Reserved Instance (RI) discount and Saving Plans discount sharing for any accounts in that organisation, including the payer account.
    \item This means that RIs and Saving Plans discounts aren't shared between any accounts that have sharing turned off.
    \item To share an RI or Savings Plans discount with an account, both accounts must have sharing turned on.
\end{itemize}

\subsubsection{AWS Organisation - Moving Accounts}
\begin{itemize}
    \item Remove the member account from the AWS organisation
    \item Send an invite to the member account from the AWS organisation
    \item Accept the invite to the new Organisation from the member account.
%Diagram
\end{itemize}

\subsubsection{Service Control Policies (SCP)}
\begin{itemize}
    \item Define allowlist or blocklist IAM actions
    \item Applied at the OU or Account level
    \item Does not apply to the Management Account
    \item SCP is applied to all the Users and Roles in the account, including Root user
    \item The SCP does not affect Service-linked roles
    \item Service-linked roles enable other AWS services to integrate with AWS organisations and can't be restricted by SCP's
    \item SCP must have an explicit Allow from the root of each OU in the direct path to the target account (does not allow anything by default)
    \item Use cases:
    \item Restrict access to certain services (for example: can't use EMR)
    \item Enforce PCI compliance by explicitly disabling services.
\end{itemize}

\subsubsection{SCP Hierarchy}
%Insert diagram
\begin{itemize}
    \item Management Account - Can do anything (no SCP apply)
    \item Account A
    \item Can do anything
    \item EXCEPT S3 (explicit Deny from Sandbox OU)
    \item EXCEPT EC2 (explicit deny)
    \item Account B and C
    \item Can do anything
    \item EXCEPT S3 (explicit Deny from Sandbox OU)
    \item Account D
    \item Can access EC2
    \item Prod OU and Account E and F
\end{itemize}

\subsubsection{SCP Examples - Blocklist and Allowlist strategies}
%Insert JSON examples

\subsubsection{IAM Policy Evaluation Logic}
%Diagram and reference to amazon docs

\subsubsection{Restricting Tags with IAM policies}
\begin{itemize}
    \item You can restrict specific tags on AWS resources
    \item Using the aws:TagKeys Condition Key
    \item Validate the Tag Keys attached to a resource against the Tag Keys in the IAM Policy
    \item Example: Allow IAM users to create EBS Volumes only if it has the "Env" and "CostCenter" Tags
    \item Use either ForAllValues (must have all keys) or ForAnyValue (must have any of these keys at a minimum)
\end{itemize}

%Add JSON tags here

\subsubsection{Using SCP to restrict creating resources without appropriate tags}
- Prevent IAM Users/Roles in the affected member accounts from creating resources if they don't have a specific Tag
% Add JSON example here

\subsubsection{AWS Organisations - Tag Policies}
\begin{itemize}
    \item Helps you standardise tags across resources in an AWS organisation
    \item Ensure consistent tags, audit tagged resources, maintain proper resources categorisation
    \item You define Tag keys and their allowed values
    \item Helps with AWS Cost Allocation Tags and Attribute-based Access Control
    \item Prevent any non-compliant tagging operations on specified services and resources
    \item Generate a report that lists all tagged/ non compliant resources
    \item Use Amazon EventBridge to monitor non-compliant tags
\end{itemize}

\subsubsection{AWS Organisation - AI Service Opt-out Policies}
\begin{itemize}
    \item Certain AWS AI services may use your content for continuous improvement of Amazon AI/ML services
    \item Example: Amazon Lex, Amazon Comprehend, Amazon Polly.........
    \item You can opt-out of having your content stored or used by AWS AI services
    \item Create an Opt-out Policy that enforces this setting across all Member accounts and AWS Regions
    \item You can opt-out all AI services or selected services
    \item Can be attached to Organisation Root, specific OU or individual Member account
\end{itemize}

\subsubsection{AWS Organisations - Backup policies}
\begin{itemize}
    \item AWS Backup enables you to create Backup Plans that define how to backup your AWS resources
    \item JSON Documents that define backup plans across an AWS Organisation
    \item Gives you granular control over backup up your resources (e.g backup frequency, time window, backup region,.....)
    \item Can be attached to Organisation Root, specific OU or individual Member account
    \item Immutable backup plans appear in Member accounts (view only)
\end{itemize}

\subsubsection{Using SCP to Deny a Region aws:ReqeustRegion}
%insert JSON example

\subsection{AWS IAM Identity Center}

\subsubsection{AWS IAM Identity Center -successor to AWS Single Sign-On- }

\begin{itemize}
    \item One login (single sign-on) for all your
    \item AWS accounts in AWS organisations
    \item Business cloud applications (e.g, Salesforce, Box, Microsoft 365)
    \item SAML2.0 enabled applications
    \item EC2 windows instances

    \item Identity Providers
    \begin{itemize}
        \item Built-in identity store in IAM identity center
        \item 3rd party Active Directory (AD), OneLogin, Okta
    \end{itemize}

    \item AWS IAM Identity Center - Login Flow
\end{itemize}

\subsubsection{AWS IAM Identity Center}
%Insert diagram

\subsubsection{AWS IAM Identity Center - Fine-grained Permissions and Assignments}
\begin{itemize}
    \item Multi-Account Permissions
    \item Manage access across AWS accounts in your AWS Organisation
    \item Permission Sets - a collection of one or more IAM policies assigned to users and groups to define AWS access
%diagrams
    \item Application Assignments
    \item SSO access to many SAML 2.0 business applications (Salesforce, Box, Microsoft 365)
    \item Provide required URLs, certificates and metadata
    \item Attribute-Based Access Control (ABAC)
    \item Fine-grained permissions based on users' attributes stored in IAM identity center identity store
    \item Example: Cost center, title, locale
    \item Use case: Define permission once, then modify AWS access by changing the attributes
\end{itemize}

\subsubsection{AWS Control Tower}
\begin{itemize}
    \item Easy way to setup and govern a secure and compliant multi-account AWS environment based on best practices
    \item Benefits:
    \begin{itemize}
        \item Automate the set up of your environment in a few clicks
        \item Automate ongoing policy management using guardrails
        \item Detect policy violations and remediate them
        \item Monitor compliance through an interactive dashboard
        \item AWS ControlTower runs on top AWS Organisations:
        \item It automatically sets up AWS Organisations to organise accounts and implement SCPs (Service Control Policies)
    \end{itemize}
\end{itemize}

\subsubsection{AWS Controller Tower - Account Factory}
\begin{itemize}
    \item Automates account provisioning and deployments
    \item Enables you to create pre-approved baselines and configuration options for AWS accounts in your organisation (e.g VPC default configuration, subnets, region, ...)
    \item Uses AWS service catalog to provision new AWS accounts
%diagram
\end{itemize}

\subsubsection{AWS Control Tower - Detect and Remediate Policy Violations}
\begin{itemize}
    \item Guardrail
    \item Provide ongoing governance for your Control Tower environment (AWS Accounts)
    \item Preventive - using SCPs (e.g Disallow Creation of Access Keys for the Root User)
    \item Detective - users AWS Config (e.g Detect Whether MFA for the Root User is Enabled)
    \item Example: identify non-compliant resources (e.g, untagged resources)
%diagram
\end{itemize}

\subsubsection{AWS Control Tower - Guardrails Levels}
\begin{itemize}
    \item Mandatory
    \item Automatically enabled and enforced by AWS control tower
    \item Example: Disallow public Read access to the Log Archive account
    \item Strongly Recommended
    \item Based on AWS best practices (optional)
    \item Example: Enable encryption for EBS volumes attached to EC2 instances
    \item Elective
    \item Commonly used by enterprises (optional)
    \item Examples: Disallow delete actions without MFA in S3 buckets
\end{itemize}


\section{AWS Resource Access Manager (RAM)}
\begin{itemize}
    \item Share AWS resources that you own with other AWS accounts
    \item Share with any account or within your Organisation
    \item Avoid resource duplication!
    \item VPC Subnets
    \item Allow to have all the resources launched in the same subnets
    \item Must be from the same AWS organisations
    \item Cannot share security groups and defaultVPC
    \item Participants can manage their own resources in there
    \item Participants can't view, modify, delete resources that belong to other participants or the owner
    \item AWS Transit Gateway
    \item Route 53 (Resolver Rules, DNS Firewall Rule Groups)
    \item License Manager Configurations
\end{itemize}

\subsubsection{AWS Resource Access Manager (RAM)}
\begin{itemize}
    \item Aurora DB Clusters
    \item ACM Private Certificate Authority
    \item CodeBuild Project
    \item EC2 (Dedicated Hosts, Capacity Reservation)
    \item AWS Glue (Catalog, Database, Table)
    \item AWS Network Firewall Policies
    \item AWS Resources groups
    \item Systems Manager Incident Manager (Contacts, Response Plans)
    \item AWS Outposts (Outpost, Site)
\end{itemize}

\subsubsection{Resource Access Manager - VPC example}
%Insert diagram here
\begin{itemize}
    \item Each account......
    \item Is responsible for it own resources
    \item Cannot view modify or delete other resources in other accounts
    \item Network is shared so.....
    \item Anything deployed in the VPC can talk to other resources in the VPC
    \item Applications are accessed easily across accounts, using a private IP
    \item Security groups from other accounts can be referenced for maximum security
    \item Use cases
    \item Applications within the same trust boundaries
    \item Applications with a high degree of interconnectivity
\end{itemize}

\subsubsection{Resource Access Manager Managed Prefix List}
\begin{itemize}
    \item A set of one or more CIDR blocks
    \item Makes it easier to configure and maintain Security Groups and Route Tables
    \item Customer-Managed Prefix List
    \item Set of CIDRs that you define and manage by you
    \item Can be shared with other AWS accounts or AWS Organisation
    \item Modify to update many security groups at once
    \item AWS-Managed Prefix List
    \item Set of CIDRs for AWS services
    \item You can't create, modify, share or delete them.
\end{itemize}

\subsubsection{Resource Access Manger Route 53 Outbound Resolver}
- Helps you scale forwarding rules to your DNS in case you have multiple accounts and VPC
%diagram


\section{Summary of Identity and Federation}

\begin{itemize}
    \item Users and Accounts all in AWS
    \item AWS Organisations
    \item AWS Control Tower to setup secure and compliant multi-account AWS environment (best practices)
    \item Federation with SAML
    \item Federation without SAML with a custom IdP (GetFederationToken)
    \item IAM Identity Center to connect to multiple AWS Accounts (Organisation) and SAML apps
    \item Web Identity Federation (not recommended)
    \item Cognito for most web and mobile applications (has anonymous mode, MFA)
    \item AWS Directory Service:
    \begin{itemize}
        \item Managed Microsoft AD - standalone or setup trust AD with on-premises, has MFA, seamless joins, RDS integration
        \item AD Connector - proxy requests to on-premises
        \item Simple AD - standalone and cheap AD-compatible with no MFA, no advanced capabilities
    \end{itemize}
    \item AWS RAM to share resource (example VPC subnets)
\end{itemize}